\begin{definition}[The network]
  \label{def:aba-network}
  \lean{PLTS.ABA.ABDY.NetworkState}\lean{PLTS.ABA.ABDY.NetworkState.corrupt}\lean{PLTS.ABA.ABDY.abdyGhostStep}\lean{PLTS.ABA.Implementation.NetworkStep.announceBound}\lean{PLTS.ABA.ABDY.NetworkStep}\lean{PLTS.ABA.ABDY.network}\leanok
  \uses{def:aba-programs, def:aba-labels, def:aba-parameters}
  The network $\mathrm{network}(P)$ holds what a process may not see: the round-tagged sent sets $sent(r,j)$ of the round messages $j$ has multicast in round $r$ (D5), the sent sets $decidedSent(j)$ of the DECIDED payloads $j$ has multicast (D12$'$), the corrupted set $F$ with its budget, the guard $k \notin F \wedge |F| < f$ on the $\mathrm{fail}$ transition (D1), and the rounds a correct process has started (D37).
  It supplies the other half of every send and every delivery, the delivery guards $m \in sent(r,j)$ and $b \in decidedSent(j)$ being the soundness of the authenticated network.
  It outputs a corrupted process's graded-agreement call and coin call, $\mathrm{byzantineCallG}$ and $\mathrm{byzantineCallW}$, under $k \in F$ and $r \in \mathit{startedRounds}$ (D11, D37), and its two injections $\mathrm{byzantineGBCA}$ and $\mathrm{byzantineDecided}$, by which a corrupted sender adds an arbitrary round message or either DECIDED bit to its own sent set and so equivocates, are its silent transitions, the round injection under the same two guards.
  It takes every input on one unguarded transition: the graded-agreement call records the called round among the started rounds at a caller outside $F$, a send records the message, and the ABA return, the replaced program's at either bit included, moves nothing (D23).
  It is Definition~\ref{def:aba-implementation} at $G = \mathbb{B} \cup \{\bot\}$, so its ghost for a round is one bit, the round's bound bit (D29, D30).
  The write $\mathrm{abdyGhostStep}$ is the round's first return $\mathrm{retG}(r,id,out)$ recording $\mathrm{boundOf}$ of the round's sent sets, the corrupted set and the outcome $out$, a later return leaving the recorded bit, and every other transition leaving the ghost alone, so the ghost is write-once.
  A replaced program takes no transition on $\mathrm{retG}$, so the round's first return is a return at an unreplaced program.
  The read of the ghost is the identity, so the announcement \leandeclnamed{announceBound}{PLTS.ABA.Implementation.NetworkStep.announceBound} fires on $\mathrm{boundG}(r,b)$ exactly when the ghost of round $r$ holds $b$, from the round's first return on, and may repeat.
  It is the bit a round's own network state holds, computed from the network's round-tagged sent sets.
  The full list of transitions: \leandecl{PLTS.ABA.ABDY.network}.
\end{definition}
